\documentclass[11pt,a4paper]{article}
\usepackage[margin=27mm]{geometry}
\usepackage[T1]{fontenc}
\usepackage{lmodern,amsmath,amssymb,amsthm,mathtools,booktabs}
\usepackage[colorlinks=true,linkcolor=blue,urlcolor=blue,citecolor=blue]{hyperref}
\usepackage{enumitem}
\usepackage{microtype}
\allowdisplaybreaks
\newtheorem{theorem}{Theorem}[section]
\newtheorem{lemma}[theorem]{Lemma}
\newtheorem{proposition}[theorem]{Proposition}
\newtheorem{corollary}[theorem]{Corollary}
\theoremstyle{definition}
\newtheorem{definition}[theorem]{Definition}
\newtheorem{remark}[theorem]{Remark}
\newcommand{\Z}{\mathbb Z}
\newcommand{\Q}{\mathbb Q}
\newcommand{\pp}{\mathcal P}
\newcommand{\NN}{\mathcal N}
\renewcommand{\AA}{\mathcal A}
\newcommand{\DD}{\mathcal D}
\newcommand{\BB}{\mathcal B}
\newcommand{\sgn}{\operatorname{sgn}}
\newcommand{\red}{\operatorname{red}}
\newcommand{\normcoord}{\operatorname{norm}}
\title{Descent and finite congruences\protect\\for \(x^2+1=2y^4\)}
\author{A self-contained exposition of the assembled argument}
\date{21 September 2026}

\begin{document}
\maketitle
\begin{abstract}
We give an elementary route to the assertion that the only positive
integer solutions of \(x^2+1=2y^4\) are \((1,1)\) and \((239,13)\).
A descent first covers a larger family of quartic equations. Explicit
rational identities then index that family by an integer recurrence.
Finally, congruences modulo suitably chosen finite powers of \(3\)
exclude every unwanted index. The prerequisites are coprime
factorization, primitive Pythagorean triples, B\'ezout's identity, and
finite polynomial algebra. Rational points on a cubic provide useful
coordinates, but no theorem about their group, rank, or generators is
assumed. No approximation estimate, reciprocity law, or infinite
\(3\)-adic construction is used.
\end{abstract}
\tableofcontents

\section{The problem and the three jobs a proof must do}

\begin{theorem}\label{thm:main}
For positive integers \(x,y\),
\[
 x^2+1=2y^4
 \quad\Longleftrightarrow\quad
 (x,y)=(1,1)\ \text{or}\ (239,13).
\]
\end{theorem}

The two displayed pairs are easily checked. The difficulty is proving
that there are no others. Factoring or applying Pythagorean
parametrization repeatedly suggests a descent, but the smaller object
usually ceases to have a coordinate equal to \(1\). Rather than discard
it, we enlarge the family to
\begin{equation}\label{eq:node}
 r^4+q^2=2s^4,\qquad r,s>0,\qquad \gcd(r,s)=1,
 \qquad q\equiv1\pmod4.
\end{equation}
We call such a triple a \emph{node}. The sign of \(q\) is allowed to be
negative. An original solution gives the node \((1,y,\pm x)\), with
the sign chosen to satisfy the last congruence.

There are three separate tasks.
\begin{enumerate}[label=\arabic*.]
\item Prove an exhaustive descent on all nodes, with a strictly
      decreasing positive integer height.
\item Put every node on an explicitly indexed sequence, keeping the
      integer numerators reduced. Equality of rational coordinates alone
      would not be enough.
\item Detect which sequence entries have the required numerator \(-1\).
      This will be a congruence argument valid at every integer index.
\end{enumerate}

The descent uses height \(s\); it does \emph{not} preserve \(r=1\).
The last task supplies exactly the information that this descent loses.
The proof is longer than a typical olympiad solution because it proves
the indexing and congruence machinery as well as the descent.

\paragraph{Historical position.}
The signed descent used here is the construction attributed to Conrad
in the working notes listed at the end. Elkies's 2014 online answer
suggested combining the cubic \(Y^2=X^3-2X\) with an obstruction modulo
powers of \(3\). We credit that proposed direction; his answer did not
provide the complete elementary argument presented here.
This manuscript supplies the exhaustive signed descent, the explicit
integer recurrence and its coordinate correspondence, and the finite
congruence proof excluding every unwanted index. Published antecedents
for the recurrence invariant and difference-operator construction are
identified in Appendix~\ref{app:scope}; their required special cases
are proved here. Priority for the complete elementary argument has not
been established. The decisive obstruction is a congruence in the
sequence index, distinct from the compared approximation and reciprocity
arguments.

\section{The exhaustive signed descent}\label{sec:descent}

We first establish the elementary arithmetic that makes the descent
complete. It is important to retain signs: neither a Pythagorean leg
nor an intermediate parameter has to be positive merely because the
original unknowns are positive.

\subsection{Two elementary parametrizations}

We use the familiar coprime-product principle: if coprime positive
integers have square product, each is a square. The same reasoning
applies prime by prime when their product is twice a square.

If \(a^2+b^2=h^2\) is primitive, \(h>0\), \(a\) is odd and
\(b\ne0\) is even, then
\[
 a=u^2-v^2,\qquad b=2uv,\qquad h=u^2+v^2,
 \qquad \gcd(u,v)=1,
\]
where \(u>0\), \(v\) may be negative, and \(u,v\) have opposite
parity. Indeed, \((h+a)/2\) and \((h-a)/2\) are positive coprime
integers with product \((b/2)^2\); take their square roots, giving
the second root the sign of \(b\). This proof allows \(a<0\).
Simultaneously changing the signs of \(u,v\) preserves the formulas.

We will also use the following equally short parametrization. If
\(a,c>0\), \(a,c\) are odd, \(b\ne0\) is even,
\(\gcd(a,b)=1\), and \(a^2+2b^2=c^2\), then
\begin{equation}\label{eq:normtwo}
 a=|t^2-2w^2|,\qquad b=2tw,\qquad c=t^2+2w^2,
 \qquad t>0,\quad \gcd(t,w)=1.
\end{equation}
To see this, factor \((c-a)/2\) and \((c+a)/2\). They are positive
and coprime, their product is \(b^2/2\), and one is odd. Thus the
odd factor is \(t^2\), the even one is \(2w^2\). Choose the sign of
\(w\) to match \(b\). In particular, \(t\) is odd; if \(4\mid b\),
then \(w\) is even.

\subsection{Constructing a smaller node}

\begin{proposition}[Signed descent]\label{prop:descent}
Every node \((x,y,z)\) with \(y>1\) is obtained from a node
\((r,s,q)\) with \(s<y\) as follows. There are integers
\(\lambda>0\), \(\mu\ne0\), and \(\delta\in\{1,-1\}\) such that
\begin{equation}\label{eq:cell}
 (r^2+2s^2)\mu=\lambda(rs+\delta q).
\end{equation}
Put
\begin{align}\label{eq:ascent}
 F&=r^2\lambda^2-2s^2\mu^2,&
 H_c&=s^2\lambda^2+r^2\mu^2,\notag\\
 U&=s^2\lambda^2-r^2\mu^2,&
 V&=2rs\lambda\mu,\notag\\
 Z_c&=U^2-V^2-2UV.
\end{align}
Then \((x,y,z)=(|F|,H_c,Z_c)\). The only node with \(y=1\)
is \((1,1,1)\).
\end{proposition}

\begin{proof}
The node equation and coprimality imply that \(x,y,z\) are odd and
pairwise coprime. For example, if \(y\) were even, \(x\) would be
odd and reduction modulo \(4\) would be impossible. An odd prime
dividing any two coordinates divides the third. Moreover,
\(z^2\equiv1\pmod{16}\), so the convention \(z\equiv1\pmod4\)
gives \(z\equiv1\pmod8\).

Set
\[
 A_0=\frac{x^2+z}{2},\qquad B_0=\frac{x^2-z}{2}.
\]
These signed integers satisfy
\[
 A_0^2+B_0^2=y^4,\qquad \gcd(A_0,B_0)=1,
 \qquad A_0\equiv1\pmod4,\quad 4\mid B_0.
\]
If either leg is zero, the node is the root; hence both are nonzero
when \(y>1\). Apply the signed Pythagorean parametrization:
\begin{equation}\label{eq:firstpyth}
 A_0=u^2-v^2,\quad B_0=2uv,\quad y^2=u^2+v^2.
\end{equation}
Here \(u\) is odd, \(v\) even, and \(\gcd(u,v)=1\).
The last equation modulo \(8\) gives \(4\mid v\).
Also
\begin{equation}\label{eq:secondnorm}
 x^2+2v^2=(u+v)^2.
\end{equation}
In particular \(u+v\ne0\). Negate both \(u,v\), if necessary, so
that \(u+v>0\).

Parametrize the last equation of \eqref{eq:firstpyth} once more:
\begin{equation}\label{eq:secondpyth}
 u=k^2-\ell^2,\quad v=2k\ell,\quad y=k^2+\ell^2,
 \qquad k>0,\quad \ell\ne0,\quad \gcd(k,\ell)=1.
\end{equation}
At this stage we assume only that \(k,\ell\) have opposite parity.
On the other hand \(\gcd(x,v)=1\), since a prime dividing both
would divide \(u\) by \eqref{eq:secondnorm}. Applying
\eqref{eq:normtwo} gives
\begin{equation}\label{eq:thirdparam}
 x=|t^2-2w^2|,\quad v=2tw,\quad u+v=t^2+2w^2,
 \qquad t>0\text{ odd},\quad w\text{ even},\quad \gcd(t,w)=1.
\end{equation}
Since \(4\mid v\), the last equality gives \(u\equiv1\pmod4\).
Thus \eqref{eq:secondpyth} forces \(k\) odd and \(\ell\) even.

Now \(k\ell=tw\). This is where the four-number decomposition
(Vierzahlensatz) enters. Set
\[
 \lambda=\gcd(k,t)>0,\qquad k=s\lambda,\qquad t=r\lambda.
\]
Then \(\gcd(r,s)=1\), and Euclid's lemma gives
\[
 \ell=r\mu,\qquad w=s\mu.
\]
Here \(r,s,\lambda\) are positive odd integers, \(\mu\ne0\) is
even, and \(\gcd(\lambda,\mu)=1\). In fact all four factors are
pairwise coprime: this follows by inserting the decompositions in
\(\gcd(k,\ell)=\gcd(t,w)=1\), together with \(\gcd(r,s)=1\).
The weaker two coprimalities just stated already suffice below.

Equating the two expressions for \(u+v\) gives
\begin{equation}\label{eq:compatibility}
 D\mu^2-2rs\lambda\mu+(r^2-s^2)\lambda^2=0,
 \qquad D=r^2+2s^2.
\end{equation}
Since \(\gcd(\lambda,\mu)=1\), reduction modulo \(\lambda\)
shows that \(\lambda\mid D\). The identity
\begin{align*}
 &(D\mu-rs\lambda)^2+\lambda^2(r^4-2s^4)\\
 &\hspace{12mm}=D\bigl(D\mu^2-2rs\lambda\mu+(r^2-s^2)\lambda^2\bigr)
\end{align*}
therefore shows that the integer
\[
 q_0=(D/\lambda)\mu-rs
\]
satisfies \(q_0^2+r^4=2s^4\). It is odd. Choose
\(q\in\{q_0,-q_0\}\) with \(q\equiv1\pmod4\), and write
\(q_0=\delta q\). This proves both the new node condition and
\eqref{eq:cell}.

The height decreases for the very simple reason
\[
 s\le k<k^2+\ell^2=y.
\]
The second inequality uses \(k\ge1\) and \(\ell\ne0\), not any
ordering between \(r\) and \(s\). Substitution into
\eqref{eq:secondpyth}--\eqref{eq:thirdparam} gives exactly
\eqref{eq:ascent}; also \(z=A_0-B_0=u^2-v^2-2uv\).
Finally, \(y=1\) forces \(x=1\), \(z=\pm1\); the sign convention
selects \(z=1\).
\end{proof}

Repeated descent terminates because its height is a positive integer.
Reversing the actual steps constructs every node from the root.
We do not need to assert that arbitrary choices of ascent parameters
always give primitive nodes: only the parameters furnished by the
proved descent will be used.

\section{An integer recurrence with reduced numerators}\label{sec:somos}

The next construction may initially look unexpected. Its purpose is
to replace a branching tree by one signed integer index. The
coordinate calculation in Section~\ref{sec:coordinates} explains the
constants. All required arithmetic properties are proved here; no
general theorem about Somos sequences is needed.

\begin{proposition}\label{prop:b}
There is a doubly infinite sequence of nonzero integers \(b_l\),
\(l\in\Z\), with
\begin{equation}\label{eq:somos}
 b_{l+2}b_{l-2}=1764b_{l+1}b_{l-1}-3107b_l^2,
 \qquad (b_{-2},b_{-1},b_0,b_1)=(13,1,1,13).
\end{equation}
It satisfies
\begin{gather}\label{eq:b-properties}
 b_{-l-1}=b_l,\qquad b_l\equiv1\pmod3,\qquad
 \gcd(b_l,b_{l+1})=\gcd(b_l,b_{l+2})=1,\\
 (b_{-2},b_{-1},b_0,b_1,b_2,\ldots)
       \equiv(5,1,1,5,5,1,1,5,\ldots)\pmod8.\notag
\end{gather}
Every term is coprime to \(1764=2^2 3^2 7^2\).
\end{proposition}

\begin{proof}
For a state of four consecutive terms define
\[
 I(a,b,c,d)=a^2d^2+1764(ac^3+b^3d)-3107b^2c^2-422abcd.
\]
The initial state has \(I=0\). Given a nonzero rational \(a\), put
\(e=(1764bd-3107c^2)/a\). Multiplication gives
\begin{equation}\label{eq:invariant}
 I(b,c,d,e)=\frac ea I(a,b,c,d).
\end{equation}
When \(I=0\), it also gives the two useful formulas
\begin{equation}\label{eq:two-e}
 ae=1764bd-3107c^2,\qquad
 b^2e=422bcd-ad^2-1764c^3.
\end{equation}

Suppose the state entries are nonzero integers, each coprime to
\(1764\), with adjacent and distance-two entries coprime. Define
\(e\) rationally as above. Both \(ae\) and \(b^2e\) are integers.
Since \(\gcd(a,b^2)=1\), a B\'ezout combination proves that \(e\)
is an integer. This step does not assume the desired divisibility
by \(a\).

For \(p=2,3,7\), the first equation in \eqref{eq:two-e} reduces to
\(ae\equiv c^2\pmod p\). Thus \(e\) is coprime to these three
primes and in particular is nonzero. If a prime divided both \(e,d\),
the preceding coprimalities would give \(p\nmid b,c\). The two
equations in \eqref{eq:two-e} would force respectively
\(p\mid3107\) and \(p\mid1764\), impossible. If a prime divided
both \(e,c\), the first equation would force \(p\mid1764\),
contrary to the coprimality just proved. These are exactly the new
adjacent and distance-two gcd conditions. The induction therefore
continues. Notice that all four entries need not be pairwise coprime:
the initial first and fourth entries are both \(13\).

The reverse recurrence is the same rule in reversed order. The
symmetric four seeds consequently give the unique backward sequence
and \(b_{-l-1}=b_l\). Alternatively, define the backward sequence by
this reflection and check the recurrence at the two overlapping seeds.

Modulo \(3\), four entries equal to \(1\) produce another \(1\).
Modulo \(8\), all entries are odd, so
\(ae\equiv4bd+5c^2\equiv1\pmod8\), provided the four-term pattern
displayed in the statement holds. It follows that \(e\equiv a\)
because each odd residue is its own inverse modulo \(8\). This
proves the pattern by induction in both directions.
\end{proof}

\begin{definition}\label{def:f}
For every integer \(l\), put
\[
 f_l=\frac{9b_l^2-b_{l-1}b_{l+1}}4,
 \qquad X_l=\frac{f_l}{b_l^2}.
\]
\end{definition}

\begin{lemma}[Primitive numerator]\label{lem:primitive}
The number \(f_l\) is an odd integer and \(\gcd(f_l,b_l^2)=1\).
If \(s>0\) is an integer and \(X_l=-1/s^2\), then
\[
 f_l=-1,\qquad |b_l|=s.
\]
Moreover \(f_{-l-1}=f_l\).
\end{lemma}
\begin{proof}
The residue pattern modulo \(8\) gives
\(b_{l-1}b_{l+1}\equiv5\), whereas \(9b_l^2\equiv1\).
Thus the numerator defining \(f_l\) is \(4\pmod8\). A prime
dividing both \(f_l\) and \(b_l\) would divide
\(b_{l-1}b_{l+1}\), contradicting adjacent coprimality. If
\(f_ls^2=-b_l^2\), then \(|f_l|\mid b_l^2\), so coprimality
forces \(|f_l|=1\). Its sign is negative, and \(s=|b_l|\) follows.
Reflection follows from that of \(b_l\).
\end{proof}

\section{Why the recurrence indexes the descent}\label{sec:coordinates}

\subsection{A coordinate change and two explicit maps}

Associate to a node the rational pair
\begin{equation}\label{eq:nodepoint}
 E(r,s,q)=\left(-\frac{r^2}{s^2},\frac{rq}{s^3}\right).
\end{equation}
The node equation becomes
\begin{equation}\label{eq:curve}
 Y^2=X^3-2X.
\end{equation}
This is a change of coordinates, not an invocation of a theorem about
elliptic curves. Its advantage is that the complicated ascent formulas
become simple compositions of fixed rational maps.

Here are those maps. Let \(P=(-1,1)\), \(Q=(9/4,-21/8)\), and
\(\NN(X,Y)=(X,-Y)\). Define
\begin{equation}\label{eq:Amap}
 m=\frac{Y+21/8}{X-9/4},\qquad
 \AA(X,Y)=(m^2-X-9/4,\ m(X-(m^2-X-9/4))-Y).
\end{equation}
For \(\DD\), first put
\begin{equation}\label{eq:Dmap}
 d=\frac{3X^2-2}{2Y},\quad X_2=d^2-2X,\quad Y_2=d(X-X_2)-Y,
 \quad e=\frac{Y_2-1}{X_2+1},
\end{equation}
and define
\[
 \DD(X,Y)=(e^2-X_2+1,\ e(X_2-(e^2-X_2+1))-Y_2).
\]
Geometrically, these formulas come from the third intersection of a
line with the cubic, followed by reflection. The first slope joins
\((X,Y)\) to \(Q\); the latter two slopes are a tangent and then a
line to \(P\). This explains the formulas but supplies no group-law
assumption. Substitution into the cubic proves the assertions below.

\subsection{Checking every denominator}

A rational number will be called \emph{regular modulo \(3\)} if its
reduced denominator is not divisible by \(3\); it then has a residue
modulo \(3\) in the usual way. Let \(\mathcal L\) consist of rational
pairs satisfying \eqref{eq:curve}, with \(X\) regular modulo \(3\)
and \(X\equiv-1\pmod3\). The equation implies that \(Y\) is also
regular modulo \(3\): a factor \(3\) in its reduced denominator
could not disappear on squaring. It also gives \(Y^2\equiv1\pmod3\).

Every node lies in \(\mathcal L\). Indeed, reduction of
\(r^4+q^2=2s^4\) modulo \(3\), with \(\gcd(r,s)=1\), shows that
none of \(r,s,q\) is divisible by \(3\). Thus
\(-r^2/s^2\equiv-1\pmod3\).

The denominator \(4X-9\) in \(\AA\) is \(2\pmod3\).
In \(\DD\), \(2Y\ne0\); moreover
\begin{equation}\label{eq:double-denom}
 X_2+1=\frac{X^4+4X^3+4X^2-8X+4}{4Y^2},
\end{equation}
whose numerator is \(1\pmod3\). All these denominators are therefore
nonzero. Direct reduction of the formulas gives
\[
 \AA(X,Y)\equiv(-1,-Y),\qquad
 \DD(X,Y)\equiv(-1,-1)\pmod3.
\]
Consequently \(\AA,\DD,\NN\) map \(\mathcal L\) into itself,
so the same denominator checks apply to every composition.

\begin{lemma}[Fixed rational identities]\label{lem:maps}
On \(\mathcal L\), \(\AA\) is a bijection with inverse
\(\NN\AA\NN\), and
\begin{equation}\label{eq:maplaws}
 \DD\AA=\AA^2\DD,\qquad \DD(P)=\AA(P),\qquad \AA(\NN P)=P.
\end{equation}
If a child is reconstructed from a parent by
\eqref{eq:cell}--\eqref{eq:ascent}, then its pair
\eqref{eq:nodepoint} is
\begin{equation}\label{eq:childD}
 \NN^{\epsilon}\DD\NN^{\eta}E(r,s,q)
 \qquad\text{for some }\epsilon,\eta\in\{0,1\}.
\end{equation}
\end{lemma}

\begin{proof}
These are finite rational identities, with denominators already
checked. Appendix~\ref{app:maps} supplies polynomial versions and a
coefficient-reduction procedure for \eqref{eq:maplaws}. In particular
\(\NN\AA\NN\AA=\mathrm{id}\); conjugating by \(\NN\) gives the
other inverse identity. The two base identities are small substitutions;
explicitly, \(\DD(P)=\AA(P)=(-1/169,239/2197)\).
For the last one, the line joining \((-1,-1)\) to \(Q\) has slope
\(-1/2\), and \eqref{eq:Amap} gives \((-1,1)\).

Here is a more compact check of the ascent identity. Write
\[
 t=r/s,\quad w=q/s^2,\quad h=(t+\delta w)/(t^2+2),
 \quad R=t^2-2h^2,\quad H=1+t^2h^2,
\]
\[
 U_0=1-t^2h^2,\quad V_0=2th,\quad
 Z_0=U_0^2-V_0^2-2U_0V_0.
\]
Then \(w^2+t^4=2\) and \eqref{eq:cell} gives \(h=\mu/\lambda\).
Consequently the child pair is
\[
 \left(-\frac{R^2}{H^2},\frac{|R|Z_0}{H^3}\right).
\]
Substitution in \eqref{eq:Dmap}, with input
\((-t^2,\delta tw)\), gives
\begin{equation}\label{eq:ascent-map-check}
 \DD(-t^2,\delta tw)
   =\left(-\frac{R^2}{H^2},-\frac{RZ_0}{H^3}\right).
\end{equation}
To check this equality, replace \(w^2\) by \(2-t^4\) after clearing
denominators; both differences are zero polynomials. The first
coordinate is unchanged by \(\NN\), and the second requires only
the sign \(-\sgn(R)\). Here \(R\ne0\): its numerator
\(F\) in \eqref{eq:ascent} is odd. Also \(H>0\).
This proves \eqref{eq:childD} with all signs retained.
\end{proof}

\subsection{The exact rational pair supplied by the sequence}

Define
\begin{equation}\label{eq:q-and-point}
 q_l=\frac{b_{l-1}b_{l+1}}{b_l^2},\qquad
 Y_l=\frac{q_l^2(q_{l-1}-q_{l+1})}{672},\qquad
 \pp_l=(X_l,Y_l).
\end{equation}
The new symbol \(q_l\) denotes a rational sequence and is distinct
from a node's signed integer coordinate \(q\).

\begin{lemma}\label{lem:actualpoint}
For every integer \(l\), the point \(\pp_l\) lies in
\(\mathcal L\), and
\[
 \pp_0=P,\qquad \pp_{l+1}=\AA(\pp_l),\qquad
 \pp_{-l-1}=\NN\pp_l.
\]
\end{lemma}
\begin{proof}
All \(q_l\) are nonzero. The recurrence and its invariant imply
\begin{align}
 q_{l-1}q_{l+1}&=\frac{1764q_l-3107}{q_l^2},\label{eq:qrt1}\\
 q_l^2q_{l+1}^2-422q_lq_{l+1}
       +1764(q_l+q_{l+1})-3107&=0.\label{eq:qrt2}
\end{align}
For example, dividing \(I(b_{l-1},b_l,b_{l+1},b_{l+2})=0\)
by \(b_l^2b_{l+1}^2\) gives \eqref{eq:qrt2}.

Here are explicit checks of both coordinates, so that a hidden
choice of sign of \(Y_l\) cannot enter. Put \(u=q_l\),
\(v=q_{l+1}\), and
\[
 J=u^2v^2-422uv+1764(u+v)-3107,\quad
 N=1764u-3107-u^2v^2.
\]
The proposed point and its successor have coordinates
\[
 X=(9-u)/4,\quad Y=N/(672v),\quad
 X'=(9-v)/4,\quad
 Y'=(u^2v^2-1764v+3107)/(672u).
\]
Direct multiplication gives
\begin{align*}
 &N^2-7056v^2\bigl((9-u)^3-32(9-u)\bigr)\\
 &\qquad=J\bigl(u^2v^2+422uv+1764u-1764v-3107\bigr).
\end{align*}
Thus \(Y^2=X^3-2X\). With
\(m=(Y+21/8)/(X-9/4)\), a second identity is
\[
 m^2-X-9/4-X'
 =\frac{J(u^2v^2+422uv+1764u+1764v-3107)}{28224u^2v^2}.
\]
Finally \(m(X-X')-Y=Y'\), an identity obtained by substituting
these four displayed fractions. Since \(J=0\), these are precisely
both formulas for \(\AA\). The denominators are legitimate:
\(u,v\ne0\) and \(X-9/4=-u/4\ne0\).

The seeds give \(b_2=1525\), \(q_0=13\), and
\(q_1=1525/169\), whence \(\pp_0=(-1,1)\).
Reflection of \(b_l\) gives reflection of \(q_l\), reverses the
difference in \eqref{eq:q-and-point}, and proves
\(\pp_{-l-1}=\NN\pp_l\).

Finally \(b_l\equiv1\pmod3\) and
\(4f_l=9b_l^2-b_{l-1}b_{l+1}\) give \(f_l\equiv-1\pmod3\).
Thus \(X_l\equiv-1\pmod3\), with denominator prime to \(3\).
The already proved curve identity now gives the required regularity
of \(Y_l\). One must use this argument rather than assume regularity
from a fraction whose displayed denominator \(672\) contains \(3\).
\end{proof}

\subsection{Turning the binary tree into a signed index}

\begin{lemma}\label{lem:binaryindex}
For every integer \(l\),
\[
 \DD(\pp_l)=\pp_{2l+1}.
\]
\end{lemma}
\begin{proof}
The assertion at \(l=0\) is \(\DD(P)=\AA(P)\).
If it holds at \(l\), then
\[
 \DD(\pp_{l+1})=\DD\AA(\pp_l)
 =\AA^2\DD(\pp_l)=\AA^2\pp_{2l+1}=\pp_{2l+3}.
\]
For the backward step, the same intertwining identity gives
\(\AA^2\DD(\pp_{l-1})=\DD(\pp_l)=\pp_{2l+1}
=\AA^2\pp_{2l-1}\). Cancel \(\AA^2\), which is injective.
Induction in both directions proves the result. No inverse of
\(\DD\), and no general group law, is used.
\end{proof}

\begin{corollary}[Exhaustive orbit coverage]\label{cor:coverage}
Every node \((r,s,q)\) satisfies
\[
 E(r,s,q)=\pp_l\qquad\text{for some }l\in\Z.
\]
\end{corollary}
\begin{proof}
Use strong induction on its positive height \(s\). The root gives
\(\pp_0\). For a nonroot, Proposition~\ref{prop:descent} gives a
smaller parent. Its point lies on the sequence by induction.
Lemma~\ref{lem:maps} expresses the child point by applying \(\DD\)
and possibly negating before and after. Lemma~\ref{lem:binaryindex}
and \(\NN\pp_l=\pp_{-l-1}\) show that all these operations stay
on the same sequence.
\end{proof}

\section{A finite congruence that excludes every other index}\label{sec:congruence}

The original condition is a condition on a reduced numerator.
For odd subscripts it has the convenient form
\begin{equation}\label{eq:G}
 G(j):=b_{2j}b_{2j-2}-9b_{2j-1}^2=-4f_{2j-1}.
\end{equation}
We prove the exact statement that is needed:
\begin{theorem}\label{thm:G}
For every integer \(j\),
\[
 G(j)=4\quad\Longleftrightarrow\quad j=0\text{ or }j=1.
\]
\end{theorem}

A bounded congruence sieve would not suffice: an arbitrarily large
index can imitate a small index modulo any one fixed modulus.
Instead, the proof chooses a finite power of \(3\) according to
the index being considered. A polynomial change that is the identity
modulo \(3\) gains another factor of \(3\) at each finite difference.
The binomial theorem then says that the first difference controls a
long jump to sufficient precision.

\subsection{Scaling before reducing}

Since \(b_l\equiv1\pmod3\), form the integers
\[
 z_l=(b_l-1)/3.
\]
This division takes place in \(\Z\), before any modular reduction.
For four consecutive scaled entries \(x,y,z,t\), define
\begin{align*}
 H(x,y,z,t)={}&1764ty+588t+x^2-2xz-2x+588y\\
             &-3107z^2-2072z-150.
\end{align*}
The elementary identity
\begin{align}\label{eq:scaledidentity}
 &1764(1+3y)(1+3t)-3107(1+3z)^2\notag\\
 &\qquad=(1+3x)\bigl(1+3(2-x+2z)\bigr)+9H(x,y,z,t)
\end{align}
turns the recurrence, after cancelling \(3\) over the integers,
into
\begin{equation}\label{eq:scaledexact}
 (1+3x)\bigl(z_{\rm new}-(2-x+2z)\bigr)=3H(x,y,z,t).
\end{equation}
Cancelling \(3\) inside \(\Z/3^N\Z\) would be invalid; this
order of operations avoids that mistake.

Fix any positive integer \(N\). In the full polynomial ring
\[
 R_N=(\Z/3^N\Z)[Z_0,Z_1,Z_2,Z_3]
\]
we work with ordinary polynomials in four variables, reducing their
coefficients modulo \(3^N\). In this notation,
the polynomial
\[
 I_N(X)=\sum_{k=0}^{N-1}(-3X)^k
\]
satisfies \((1+3X)I_N(X)=1\), since the omitted final power contains
\(3^N\). Thus the actual scaled transition modulo \(3^N\) is
represented by the polynomial map
\begin{equation}\label{eq:Vmap}
 V_N(x,y,z,t)=
 (y,z,t,\ 2-x+2z+3H(x,y,z,t)I_N(x)).
\end{equation}
No degree truncation of the variables is being made. The coefficient
ring is finite, but polynomials of every finite degree are allowed.

Modulo \(3\), this map is the affine map
\[
 L(x,y,z,t)=(y,z,t,2-x+2z).
\]
Six substitutions give, over the integers,
\[
 L^6(x,y,z,t)=(-2x+3z+6,-2y+3t+6,-3x+4z+12,-3y+4t+12).
\]
Consequently
\begin{equation}\label{eq:nearidentity}
 U_N:=V_N^6\equiv\mathrm{id}\pmod3
 \quad\text{coefficient by coefficient}.
\end{equation}

\subsection{Finite differences gain divisibility}

Let \(E\) denote substitution of \(U_N\) in a polynomial of
\(R_N\), and set \(\Delta=E-\mathrm{id}\). Here an operator is
just a rule that sends a polynomial to another polynomial; no analysis
is involved. Both rules respect addition. From \eqref{eq:nearidentity},
\[
 \Delta h\in3R_N\quad(h\in R_N),\qquad
 \Delta(3^a h)=3^a\Delta h\in3^{a+1}R_N.
\]
Membership in \(3^aR_N\) means that every coefficient is a multiple
of \(3^a\).
This follows by expanding each monomial after replacing each variable
by itself plus a polynomial divisible by \(3\).

Use the observable
\[
 g(Z_0,Z_1,Z_2,Z_3)=(1+3Z_0)(1+3Z_2)-9(1+3Z_1)^2.
\]
Its nonconstant part has a factor \(3\):
\[
 g=-8+3(Z_0+Z_2+3Z_0Z_2-18Z_1-27Z_1^2).
\]
Since \(\Delta\) kills constants, induction gives
\begin{equation}\label{eq:diffdiv}
 \Delta^r g\in3^{r+1}R_N\qquad(r\ge1).
\end{equation}
At the actual integer state
\[
 Z(j)=(z_{2j-2},z_{2j-1},z_{2j},z_{2j+1})
\]
we have \(g(Z(j))=G(j)\), and the six-step map sends
\(Z(j)\) to \(Z(j+3)\) modulo \(3^N\).

\begin{lemma}[Long-jump congruence]\label{lem:lifting}
If \(j\in\Z\), \(m>0\), \(e\ge0\), and \(3^e\mid m\), then
\begin{equation}\label{eq:lifting}
 G(j+3m)-G(j)
 \equiv m\bigl(G(j+3)-G(j)\bigr)\pmod{3^{e+3}}.
\end{equation}
Also, for \(a(j)=G(j+3)-G(j)\),
\begin{equation}\label{eq:increment}
 9\mid a(j),\qquad a(j+3)\equiv a(j)\pmod{27}.
\end{equation}
\end{lemma}
\begin{proof}
Choose the finite precision \(N=e+3\). The ordinary finite binomial
identity for the commuting operators \(\mathrm{id}\) and \(\Delta\)
gives
\[
 E^mg-g-m\Delta g=\sum_{r=2}^{m}\binom mr\Delta^r g.
\]
This identity can also be proved directly by induction on \(m\),
using Pascal's identity for the binomial coefficients.
Each term on the right vanishes modulo \(3^{e+3}\). To check this,
write \(v_3(t)\) for the exponent of \(3\) in a nonzero integer
\(t\), merely a finite prime-factor count. The identity
\[
 r\binom mr=m\binom{m-1}{r-1}
\]
implies
\[
 v_3\binom mr\ge\max\{0,e-v_3(r)\}.
\]
Combining with \eqref{eq:diffdiv}, the term is divisible by
\[
 3^{r+1+\max\{0,e-v_3(r)\}}.
\]
For every \(r\ge2\),
\(v_3(r)\le r-2\): check \(r=2,3\) directly; for \(r\ge4\),
\(3^{r-2}>r\), by induction. Thus the exponent is at least
\(e+3\). Evaluation at \(Z(j)\) proves \eqref{eq:lifting}.

The first assertion of \eqref{eq:increment} is
\eqref{eq:diffdiv} with \(r=1\) and precision \(2\).
With \(r=2\) and precision \(3\), it gives
\(27\mid G(j+6)-2G(j+3)+G(j)\), which is the second assertion.
\end{proof}

Everything here is a finite sum in a ring with a specified modulus.
There is no limiting series, no completion of \(\Q\), and no
inverse substitution operator. To handle a negative shift one simply
applies the positive-shift congruence from its earlier endpoint.

\subsection{The small residues force uniqueness}

The initial recurrence modulo \(27\) gives the short table
\[
\begin{array}{c|rrrrrrrrrrr}
 l&-2&-1&0&1&2&3&4&5&6&7&8\\\hline
 b_l\pmod{27}&13&1&1&13&13&22&22&1&1&4&4
\end{array}
\]
All divisions are legitimate modulo \(27\), since each entry is
\(1\pmod3\). Therefore
\begin{equation}\label{eq:Gtable}
\begin{array}{c|rrrrr}
 j&0&1&2&3&4\\\hline
 G(j)\pmod{27}&4&4&7&13&22
\end{array},
 \qquad G(0)=G(1)=4\ \text{exactly}.
\end{equation}
In particular \(a(0)\equiv9\) and \(a(1)\equiv18\pmod{27}\).

\begin{proof}[Proof of Theorem~\ref{thm:G}]
If \(j\equiv2\pmod3\), the first part of
\eqref{eq:increment} and finite telescoping give
\(G(j)-4\equiv3\pmod9\), so \(G(j)\ne4\).

Otherwise write \(j=c+3m\) with \(c=0\) or \(1\).
If \(m=0\), the exact seeds give \(G(j)=4\). Suppose \(m>0\)
and write \(m=3^eu\), \(3\nmid u\). The long-jump congruence gives
\[
 G(j)-4\equiv m a(c)\pmod{3^{e+3}}.
\]
Since \(a(c)\) is divisible by \(9\) and not by \(27\), the
right side is divisible by \(3^{e+2}\) but not by \(3^{e+3}\).
It is a nonzero residue, so \(G(j)\ne4\).

If \(m<0\), put \(h=-m>0\), with \(3^e\) the exact power of
\(3\) dividing \(h\). Apply \eqref{eq:lifting} from \(j\) forward
by \(3h\) to \(c\):
\[
 4-G(j)\equiv h a(j)\pmod{3^{e+3}}.
\]
The second part of \eqref{eq:increment}, telescoped between \(j\)
and \(c\), gives \(a(j)\equiv a(c)\pmod{27}\). Since
\(3^e\mid h\), it follows that
\(3^{e+3}\mid h(a(j)-a(c))\), so this replacement loses no
required precision. The same
nonzero-residue argument applies. This covers every signed index.
\end{proof}

The conclusion is universal, not a computation up to a large bound.
For each index the proof supplies one finite modulus that rules it
out. The formal obstruction used in this exposition is precisely
Theorem~\ref{thm:G}; no stronger all-index valuation formula is needed.

\section{Returning to the original equation}\label{sec:finish}

\begin{proof}[Proof of Theorem~\ref{thm:main}]
Let \(x,y>0\) satisfy \(x^2+1=2y^4\). Choose the sign of \(x\)
that is \(1\pmod4\); then \((1,y,\pm x)\) is a node.
Corollary~\ref{cor:coverage} supplies an integer \(l\) such that its
point is \(\pp_l\). In particular
\[
 \frac{f_l}{b_l^2}=-\frac1{y^2}.
\]
The primitive-numerator lemma gives \(f_l=-1\) and \(y=|b_l|\).

Exactly one of \(l\) and \(-l-1\) is odd. Reflection preserves
\(f_l\) and \(b_l\), so choose the odd index and write it as
\(l=2j-1\). Equation~\eqref{eq:G} gives \(G(j)=4\).
Theorem~\ref{thm:G} forces \(j=0\) or \(j=1\), hence
\(l=-1\) or \(l=1\). The seeds now give \(y=1\) or \(y=13\).
The original equation yields \(x=1\) or \(x=239\), respectively.
Both pairs indeed satisfy it.
\end{proof}

The descent did terminate, but it terminated in a larger family.
The reduced numerator and the finite congruence recovered the condition
that could not be preserved during descent. This is the role of the
recurrence: it carries the extra arithmetic condition through the
entire family, rather than seeking a smaller original solution at
every step.

\appendix
\section{Finite algebra certificates for the rational maps}\label{app:maps}

This appendix records explicit expressions for the map identities in
Lemma~\ref{lem:maps}. The computations are finite polynomial
identities; they do not import associativity or a theorem about
rational points. The formulas are included to make the large algebra
reproducible, rather than conceal it behind elliptic-curve notation.

\subsection{Polynomial triples and normalization}

For a triple \(C=(X,Y,Z)\) with \(Z\ne0\), write
\[
 \normcoord(C)=(X/Z^2,Y/Z^3),\qquad
 \NN C=(X,-Y,Z).
\]
The curve relation for such a triple is
\(Y^2=X^3-2XZ^4\). The polynomial triples below are defined for
every input. Assertions about their normalized rational maps are made
under the further hypothesis \(\normcoord(C)\in\mathcal L\).
On this domain the image denominators are nonzero, as proved in the
main text. For residue calculations we use representatives with
\(Z\) a unit modulo \(3\); the choice \(Z=1\) always suffices.
The following polynomial maps give the rational maps in the main text
after normalization:
\begin{align*}
 A_X(C)&=(4X+9Z^2)(9X-8Z^2)+84YZ,\\
 A_Y(C)&=-42(4X^3+27X^2Z^2-24XZ^4-18Z^6)
                  -(844X-135Z^2)YZ,\\
 A_Z(C)&=4X-9Z^2.
\end{align*}
For the other map set
\begin{align*}
 H(C)&=X^4+4X^3Z^2+4X^2Z^4-8XZ^6+4Z^8,\\
 R(C)&=X^4-12X^2Z^4+4Z^8+2(X^2+2Z^4)YZ,\\
 S(C)&=(X^2-4XZ^2+2Z^4)(X^2+4XZ^2+2Z^4)\\
     &\qquad\cdot(X^4-4X^3Z^2+4X^2Z^4+8XZ^6+4Z^8)\\
     &\quad-4(X^2+2Z^4)
           (X^4+8X^3Z^2+4X^2Z^4-16XZ^6+4Z^8)YZ,
\end{align*}
and \(D(C)=(-R(C)^2,R(C)S(C),H(C))\).
Substitution in the chord formulas yields
\[
 \normcoord(A(C))=\AA(\normcoord(C)),\qquad
 \normcoord(D(C))=\DD(\normcoord(C)).
\]
The weighted degrees make these independent of a common change
\((X,Y,Z)\mapsto(t^2X,t^3Y,tZ)\), \(t\ne0\).

For checking \(\AA^2\), an intermediate triple reduces the size
of the computation. Set \(a=12769\), \(d=84\), \(c=900271\), and
\begin{align*}
 B_X(C)&=(d^2X+aZ^2)(aX-2d^2Z^2)-2cdYZ,\\
 B_Y(C)&=cd^3X^3+3acdX^2Z^2-6cd^3XZ^4-2acdZ^6\\
       &\quad+\bigl((-a^3+6ad^4)Z^2-(3a^2d^2-2d^6)X\bigr)YZ,\\
 B_Z(C)&=d^2X-aZ^2.
\end{align*}
Then \(\normcoord(B(C))=\AA^2(\normcoord(C))\).
The denominator \(B_Z\) is nonzero on our locus, since it is
\(-Z^2\pmod3\) whenever \(Z\) is a unit modulo \(3\).
Weighted homogeneity transfers this nonvanishing from the representative
with \(Z=1\) to every representative with \(Z\ne0\).

\subsection{An explicit coefficient-reduction test}

Put \(c_0(x)=x^3-2x\). Every polynomial in \(x,y\), subject to
\(y^2=c_0(x)\), has a reduced expression \(u(x)+v(x)y\).
For a polynomial \(T=\sum t_{ij}x^iy^j\), its two coefficients
are explicitly
\begin{equation}\label{eq:reduction}
 \red(T)=\left(
 \sum_{i,k}t_{i,2k}x^ic_0(x)^k,
 \sum_{i,k}t_{i,2k+1}x^ic_0(x)^k\right).
\end{equation}
Multiplication of reduced pairs is therefore just
\[
 (u,v)(r,s)=(ur+c_0vs,\ us+vr).
\]
Here is an explicit check that the polynomial \(D\) agrees with
the tangent-and-chord definition, including its second coordinate.
For \(C=(x,y,1)\), abbreviate its \(H,R,S\) by the same letters
and put
\[
 K=x^6-10x^4-20x^2+8,\qquad T=K-8y^3.
\]
The tangent formulas reduce to
\[
 X_2=\frac{(x^2+2)^2}{4y^2},\qquad
 Y_2=\frac{K}{8y^3},\qquad
 e=\frac{T}{2yH}.
\]
Using the pair-multiplication rule above, both of the following
polynomials have reduced pair \((0,0)\):
\begin{align*}
 &T^2-(x^2+2)^2H^2+4y^2H^2+4y^2R^2,\\
 &T\bigl((x^2+2)^2H^2+4y^2R^2\bigr)-KH^3-8y^3RS.
\end{align*}
The first identity, divided by \(4y^2H^2\), gives the first
coordinate \(-R^2/H^2\); the second, divided by \(8y^3H^3\),
gives the second coordinate \(RS/H^3\). Their denominators are
nonzero on \(\mathcal L\). Thus the map in the main text is exactly
the polynomial map used in the finite certificates, not merely a map
with the same first coordinate.

This provides a short way to verify the identities without expanding
large powers of \(y\). For example, let \(C=(x,y,1)\) and
\(C'=\NN A\NN A(C)\). Direct substitution gives the exact
first-coordinate certificate
\[
 C'_X-x(C'_Z)^2=(254016-112896x)(y^2-x^3+2x),
\]
and the second-coordinate certificate is the same curve factor times
\begin{align*}
 &-37933056x^3-222641664x^2-285852672xy+827632512x\\
 &\qquad-37933056y^2+444019968y-564350976.
\end{align*}
Thus both reduced pairs are zero.

Here is the full list of polynomial checks required by the map lemma.
In this table \(C=(x,y,1)\), and every entry in the right column
is the result of the explicit reduction \eqref{eq:reduction}.
\[
\renewcommand{\arraystretch}{1.25}
\begin{array}{p{0.76\linewidth}|c}
\text{Polynomial to reduce}&\red\text{ value}\\\hline
 $A_Y(C)^2-A_X(C)^3+2A_X(C)A_Z(C)^4$ & $(0,0)$\\
 $D_Y(C)^2-D_X(C)^3+2D_X(C)D_Z(C)^4$ & $(0,0)$\\
 $A_X(A(C))B_Z(C)^2-B_X(C)A_Z(A(C))^2$ & $(0,0)$\\
 $A_Y(A(C))B_Z(C)^3-B_Y(C)A_Z(A(C))^3$ & $(0,0)$\\
 $(\NN A\NN A(C))_X-x(\NN A\NN A(C))_Z^2$ & $(0,0)$\\
 $(\NN A\NN A(C))_Y-y(\NN A\NN A(C))_Z^3$ & $(0,0)$\\
 $D_X(A(C))B_Z(D(C))^2-B_X(D(C))D_Z(A(C))^2$ & $(0,0)$\\
 $D_Y(A(C))B_Z(D(C))^3-B_Y(D(C))D_Z(A(C))^3$ & $(0,0)$
\end{array}
\]
Each row is an ordinary polynomial identity after substituting the
displayed definitions; \eqref{eq:reduction} specifies every coefficient
operation needed to check it. The first two rows preserve the curve,
the next two identify \(B\) with two applications of \(A\), the next
two prove the inverse, and the last two prove \(\DD\AA=\AA^2\DD\).
These computations involve both coordinates. Checking only the first
coordinate would leave an unacceptable sign ambiguity.

For a compact additional certificate of curve preservation under
\(A\), its first row before reduction equals
\[
 (-49712x^2-1265112x-592704y+1542321)(y^2-x^3+2x).
\]
The denominators have deliberately not been cancelled in the table.
Their nonvanishing is the separate modular argument in the main text,
so polynomial equality is never being used to justify an undefined
rational expression.

\section{Scope, dependencies, and provenance}\label{app:scope}

\paragraph{What is elementary here.}
The proof uses factorization into primes only for elementary gcd and
divisibility facts, not a classification of units in an algebraic
number field. The cubic is used through displayed rational formulas
and finite polynomial identities. The congruence proof uses the full
polynomial ring over a finite residue ring and the finite binomial
theorem. It does not use an infinite local expansion or a theorem
bounding its zeros. Calling these ingredients elementary is a
description of what has been proved, not a claim that the resulting
proof is short.

\paragraph{Formal verification.}
The complete argument has been checked in Lean~4.19.0. All 34 local
modules passed a rebuild, ending in the unconditional integer,
natural-number, and equivalence forms of the classification in
\texttt{parent\_NewElementary.lean}. This includes the signed descent,
both rational coordinates, denominator conditions, primitive numerator,
finite binomial congruence, and the all-integer exclusion theorem.

A separate recursive audit followed the types and proof terms of those
three final theorems and thirteen supplementary results. Its combined
closure contains 10,856 declarations, including 531 local declarations.
It found no proof placeholders, additional axioms, unsafe dependencies,
old approximation-proof dependencies, or prohibited reciprocity,
elliptic-curve, series, or completion dependencies. The only foundational
principles are Lean's standard propositional extensionality, classical
choice, and quotient soundness. The audit also confirms that all thirteen
required logical bridges occur in each final classification proof.

The audit source, verification ledger, per-module logs, and recorded
source and output hashes accompany the manuscript.
Formal checking establishes the stated mathematical conclusions;
historical attribution and the description of elementary prerequisites
are separate assessments, not conclusions of the Lean kernel.

The universal congruence theorem used here is the uniqueness statement
for \(G\). Numerical experiments with a stronger exact valuation law
are not premises, and no stronger formal statement is claimed.

\paragraph{Provenance and further orientation.}
The following references identify related work and motivation, with
the scope of each connection stated explicitly. They are not substitutes
for lemmas in the proof.
\begin{enumerate}[label={[\arabic*]}]
\item Noam D.~Elkies, answer to \emph{A Pell equation inside a Pell
equation}, Mathematics Stack Exchange, posted 10 May 2014; the relevant
postscript was added 11 May 2014:
\href{https://math.stackexchange.com/a/788692}{Elkies's answer on Mathematics Stack Exchange}.
This online outline suggests the cubic-and-congruence direction.
It does not establish the explicit Somos recurrence or the all-index
finite-congruence lemma used here, and is not a completed elementary
proof published in a paper or book.
\item \emph{Conrad's signed quartic tree and the compact synchronization
obstruction}, inherited project working notes, file
\texttt{conrad\_tree\_vierzahlensatz\_synchronization.md}.
The attribution to Conrad is retained from those notes; a primary
bibliographic identification remains to be supplied. The full descent
needed here has been proved in Section~\ref{sec:descent}, including
negative parameters and strict height.
\item Andrew N.~W.~Hone and Christine Swart,
\emph{Integrality and the Laurent phenomenon for Somos 4 and Somos 5
sequences}, Mathematical Proceedings of the Cambridge Philosophical
Society \textbf{145} (2008), 65--85.
\href{https://kar.kent.ac.uk/id/document/1522}{Author manuscript}.
Equation~(1.13) specializes exactly to our recurrence invariant with
parameters \(1764,-3107,422\); Theorem~1.2 gives the general connection
with cubic-curve coordinates. The necessary integrality and coordinate
identities are proved directly in this manuscript.
\item Bjorn Poonen, \emph{p-adic interpolation of iterates}, Bulletin
of the London Mathematical Society \textbf{46} (2014), 525--527.
\href{https://math.mit.edu/~poonen/papers/p-iteration.pdf}{Author version}.
The proof on page~2 uses the substitution-difference operator,
increasing coefficient divisibility, and a binomial expansion of
iterates. Our congruence argument shares this mechanism, but proves
the required finite statements without assuming the analytic
interpolation theorem or a convergence result.
\end{enumerate}

\end{document}
